\documentclass[12pt,a4paper]{article}

\usepackage[margin=1in]{geometry}
\usepackage[T1]{fontenc}
\usepackage{lmodern}
\usepackage{microtype}
\usepackage{xcolor}
\usepackage{hyperref}
\usepackage{enumitem}
\usepackage{longtable}
\usepackage{array}
\usepackage{titlesec}

\definecolor{Accent}{HTML}{4F46E5}
\definecolor{Muted}{HTML}{475569}
\definecolor{CodeBg}{HTML}{F1F5F9}

\hypersetup{
  colorlinks=true,
  linkcolor=Accent,
  urlcolor=Accent,
  citecolor=Accent,
  pdftitle={Salim Shrestha Portfolio Documentation},
  pdfauthor={Salim Shrestha}
}

\titleformat{\section}
  {\Large\bfseries\color{Accent}}
  {\thesection}
  {0.7em}
  {}

\titleformat{\subsection}
  {\large\bfseries}
  {\thesubsection}
  {0.7em}
  {}

\setlist[itemize]{leftmargin=1.5em}
\setlist[enumerate]{leftmargin=1.5em}
\renewcommand{\arraystretch}{1.25}

\newcommand{\code}[1]{\colorbox{CodeBg}{\texttt{#1}}}

\begin{document}

\begin{titlepage}
  \centering
  \vspace*{2cm}
  {\Huge\bfseries Salim Shrestha Portfolio\par}
  \vspace{0.6cm}
  {\Large React Portfolio Application Documentation\par}
  \vspace{1.2cm}
  {\large Vite + React 18 + TailwindCSS\par}
  \vfill
  \begin{tabular}{>{\bfseries}rl}
    Owner: & Salim Shrestha \\
    Role: & Web Developer / Computer Engineer \\
    Location: & Kathmandu, Nepal \\
    Repository: & \texttt{react-portfolio} \\
    Document Date: & May 26, 2026 \\
  \end{tabular}
  \vfill
\end{titlepage}

\tableofcontents
\newpage

\section{Project Overview}

The Salim Shrestha portfolio is a clean, modern single-page web application
implemented with Vite, React 18, and TailwindCSS. It presents Salim's personal
profile, technical skills, GitHub projects, GitHub public statistics, and
contact links in a responsive portfolio interface.

The app is designed as a client-side React application. It uses pre-seeded
profile information for stable personal details and fetches live public GitHub
data on first page load. If the GitHub API is unavailable, the application
falls back to static data so that the portfolio remains usable.

\section{Owner Details}

\begin{longtable}{>{\bfseries}p{0.25\linewidth}p{0.65\linewidth}}
Name & Salim Shrestha \\
Title & Web Developer / Computer Engineer \\
Location & Kathmandu, Nepal \\
Email & \href{mailto:salim9shrestha@gmail.com}{salim9shrestha@gmail.com} \\
Photo & \url{https://www.ssanu.com.np/images/salim.jpg} \\
CV & \url{https://www.ssanu.com.np/images/Salim-Shrestha-cv.pdf} \\
GitHub & \url{https://github.com/salimshre} \\
LinkedIn & \url{https://www.linkedin.com/in/salimshrestha} \\
Facebook & \url{https://www.facebook.com/salimshre} \\
Instagram & \url{https://www.instagram.com/salimsshrestha/} \\
YouTube & \url{https://www.youtube.com/c/technodynamic} \\
\end{longtable}

\section{Technology Stack}

\begin{longtable}{>{\bfseries}p{0.25\linewidth}p{0.65\linewidth}}
Frontend Library & React 18 \\
Build Tool & Vite \\
Styling & TailwindCSS \\
Language & Plain JavaScript and JSX \\
Icons & React Icons \\
Data Fetching & Native \code{fetch} inside a custom React hook \\
Package Manager & npm \\
\end{longtable}

\section{Repository Structure}

\begin{verbatim}
react-portfolio/
  index.html
  package.json
  vite.config.js
  tailwind.config.js
  postcss.config.js
  src/
    main.jsx
    App.jsx
    index.css
    components/
      Navbar.jsx
      Hero.jsx
      About.jsx
      Skills.jsx
      Projects.jsx
      GitHubStats.jsx
      Contact.jsx
      Footer.jsx
    hooks/
      useGitHub.js
  Documentation/
    portfolio_documentation.tex
\end{verbatim}

\section{Application Architecture}

The application starts at \code{src/main.jsx}, where React mounts \code{App}
into the \code{root} element defined in \code{index.html}. The \code{App}
component coordinates global page behavior, including:

\begin{itemize}
  \item loading GitHub data through \code{useGitHub};
  \item applying the persisted light or dark theme;
  \item tracking the active page section for scroll-spy navigation;
  \item rendering the page sections in the required single-page order.
\end{itemize}

\subsection{Component Responsibilities}

\begin{longtable}{>{\bfseries}p{0.28\linewidth}p{0.62\linewidth}}
\code{Navbar.jsx} & Fixed navigation bar, section links, active section highlight, and dark mode toggle. \\
\code{Hero.jsx} & Full-screen hero with name, location, animated role text, avatar, hire button, and CV link. \\
\code{About.jsx} & Short profile paragraph, professional title, graduation year context, and location badge. \\
\code{Skills.jsx} & Skill list with animated progress bars. \\
\code{Projects.jsx} & GitHub repository cards with language, stars, description, and external links. \\
\code{GitHubStats.jsx} & Public repository, follower, and following counters. \\
\code{Contact.jsx} & Email link and social profile icon links. \\
\code{Footer.jsx} & Footer text required by the project prompt. \\
\code{useGitHub.js} & Shared GitHub API fetch logic, loading state, error state, and fallback data. \\
\end{longtable}

\section{Data Flow}

The custom hook \code{useGitHub} fetches data from the public GitHub API:

\begin{itemize}
  \item \url{https://api.github.com/users/salimshre}
  \item \url{https://api.github.com/users/salimshre/repos?sort=updated&per_page=6}
\end{itemize}

The user endpoint provides avatar, biography, repository count, followers,
and following count. The repository endpoint provides the six most recently
updated repositories, including name, description, language, star count, and
GitHub URL.

The hook initializes the UI with fallback values and then replaces them with
live data when the API request succeeds. If the request fails, the app displays
the fallback profile and project data and shows an explanatory message in the
Projects section.

\section{User Interface Features}

\begin{itemize}
  \item Single-page layout with smooth scrolling.
  \item Fixed navigation with active section highlighting.
  \item Full-screen hero section with animated role typewriter text.
  \item Responsive design for mobile, tablet, and desktop screens.
  \item Dark mode toggle persisted in \code{localStorage}.
  \item Skeleton loading cards while GitHub data is being fetched.
  \item Graceful GitHub API fallback behavior.
  \item Contact section with email and social links.
\end{itemize}

\section{Build and Run Instructions}

\subsection{Install Dependencies}

\begin{verbatim}
npm install
\end{verbatim}

\subsection{Run Development Server}

\begin{verbatim}
npm run dev
\end{verbatim}

The local Vite development server is normally available at:

\begin{verbatim}
http://127.0.0.1:5173
\end{verbatim}

\subsection{Build Production Files}

\begin{verbatim}
npm run build
\end{verbatim}

The production output is generated in the \code{dist/} directory.

\subsection{Preview Production Build}

\begin{verbatim}
npm run preview
\end{verbatim}

\section{Environment Configuration}

The application can optionally use a GitHub token to raise GitHub API rate
limits. Create a local \code{.env} file and define:

\begin{verbatim}
VITE_GITHUB_TOKEN=your_github_token_here
\end{verbatim}

This file should remain local and must not be committed to version control.

\section{Blank Page Issue and Resolution}

During development, the app displayed a blank white page while the Vite server
was running. The cause was a Vite configuration setting that disabled
dependency optimization in development:

\begin{verbatim}
optimizeDeps: command === "serve" ? { disabled: "dev" } : {}
\end{verbatim}

With dependency optimization disabled, the browser received raw CommonJS React
modules such as \code{module.exports}, which are not valid in the browser's
native ES module runtime. This prevented React from mounting and resulted in a
blank page.

The fix was to remove the custom \code{optimizeDeps} override from
\code{vite.config.js}. Vite now pre-bundles React dependencies correctly and
serves browser-compatible modules from \code{/node\_modules/.vite/deps/}.

\section{Prompt Compliance Summary}

\begin{longtable}{>{\bfseries}p{0.42\linewidth}p{0.18\linewidth}p{0.30\linewidth}}
Requirement & Status & Notes \\
Vite + React 18 & Complete & Implemented with functional components. \\
TailwindCSS styling & Complete & \code{src/index.css} contains Tailwind directives only. \\
GitHub public API data & Complete & User and repository endpoints are fetched. \\
Fallback on GitHub errors & Complete & Static user and project data are used on failure. \\
Skeleton loaders & Complete & Projects and stats display loading placeholders. \\
Single-page sections & Complete & Hero, About, Skills, Projects, GitHub Stats, Contact, and Footer are present. \\
Dark mode toggle & Complete & Theme is persisted with guarded \code{localStorage} usage. \\
Smooth scroll-spy navigation & Complete & Intersection Observer tracks active sections. \\
Responsive layout & Complete & Tailwind mobile-first classes are used throughout. \\
Footer text & Complete & Footer matches the required 2025 wording. \\
README setup notes & Complete & Includes setup, build, preview, and optional token notes. \\
\end{longtable}

\section{Conclusion}

The repository implements the requested portfolio application with the expected
React component structure, GitHub API integration, fallback handling, responsive
TailwindCSS layout, dark mode, and documentation. The development server blank
page issue was resolved by restoring Vite's default dependency optimization
behavior.

\end{document}
